% Current review descendant; frozen semester PDF and absorbed archive unchanged.
\documentclass[11pt]{article}
\usepackage[margin=1in]{geometry}
\usepackage{amsmath,booktabs,url}
\title{Fraud Side Study: Evidence Boundaries and Withdrawn In-Sample Analyses}
\author{Arvind C R \and Ramesh Prakash Guledgudd}
\date{Source correction, 18 October 2026; not a new release}
\begin{document}
\maketitle
\begin{abstract}
This corrected descendant retains a bounded synthetic-tabular observation,
not a deployment or language-model sensor result. The historical downstream
calibration, feature-ablation, cost and cross-domain tables are withdrawn
from current scientific conclusions: their fit/evaluation populations and
iteration-specific definitions are not consistently reconciled. No retraining,
new model evaluation, or numerical replacement is represented as completed.
The historical semester PDF and absorbed source archive remain historical records.
\end{abstract}
\section{Dataset and observation units}
The supplied CSVs contain 50,000 synthetic rows (719 positives), a
40,000-row training split (575 positives), and a 10,000-row test split
(144 positives). They are not institutional transaction data. The generator
uses anonymous numeric features; row-wise mean, standard deviation, maximum
and minimum are deterministic functions of the existing features. They are
neither VLM perception nor LLM-produced regulatory narratives, and introduce
no independent unstructured information. Anonymous feature names also do
not establish a semantic transfer benefit from a pretrained language model.
\section{Retained scoring evidence and noncomparability}
The root \texttt{xgboost\_results.json} records ROC-AUC
0.794241240530303, precision 0.723404255319149, recall
0.2361111111111111 and F1 0.35602094240837695. Its recorded inference duration
is 0.006184816360473633 seconds per 10,000-row batch, not per-request latency.
The alternative historical quick-artifact AUC 0.7955 has a different
configuration and must not be merged into this result. Neither number
establishes production validity. Historical 0.975/0.948 figures are not
reconstructed and are excluded from current empirical conclusions.

The historical Qwen SFT record used a 500-row, 20\%-positive enriched subset,
not the tree's entire 10,000-row natural-rate evaluation population.
Its reported accuracy 0.792 is below the constant-negative baseline
$400/500=0.800$; AUC 0.48268 does not support useful ranking. This is not a
matched-row head-to-head effect estimate. The enriched subset bounds inference
work; fewer positives and fewer negatives do not improve AUC precision.
An exact common-row manifest and matched scoring are needed before a
comparative claim. Missing XGBoost accuracy on that subset is not imputed.
\section{Withdrawal of downstream conclusions}
The historical calibration pipeline reads \texttt{fraud\_data.csv} for fitting
and \texttt{test\_data.csv} for evaluation; the latter is contained in the
former. Consequently its apparent test scores are in-sample, not leakage-free
held-out evidence. Within-protocol comparisons do not cure differential
memorisation across feature sets. All downstream historical calibration,
feature-ablation, noise, threshold-transfer and cost tables are excluded
from generalisation or operational recommendations in this descendant.
Their bootstrap intervals cannot repair the fit/evaluation overlap.
A valid replacement would fit only disjoint training rows and reserve
independent calibration/evaluation units; that rerun is not performed here.

The historical hybrid modes charged a fee without consistently reranking
alerts using augmented scores; their cost disadvantage is not empirical
sensor utility. Deterministic aggregates require no per-row LLM invocation.
Any hypothetical token price is a scenario assumption, not measured cost.
No VLM extraction, compliance-text quality, analyst-time saving or 85-fold
cost reduction is validated. The four-feature study can at most motivate a
redundant-feature negative control, not sensor or scribe effectiveness.

The historical inverted alert threshold is not retained. For calibrated fraud
probability $p$, investigation cost $c$ on every alert and missed-fraud loss
$L$, compare $c$ with $pL$: alert if $p\ge c/L$. If the cost applies only to
false alerts, the different threshold is $c/(L+c)$. These conventions must not
be mixed, and empirical threshold selection must be out of sample.

Conflicting top-$K$ counts, per-cohort call totals, calibration values,
bootstrap tie rules, and cross-domain D-label definitions are not reconciled
by selecting a convenient iteration. In particular ECE 0.99 is worse, not
better, than 0.17--0.21. A deterministic identical-score comparison has
zero decision flips by construction. Unrelated RL telemetry and evolving
cross-paper density tables are excluded, not evidence about fraud decisions.
\section{Remaining requirements}
Before this study can support a comparative publication, recover the exact
row manifests, fit/configuration identities, predictions and costs; verify
train/test disjointness, then recompute matched-row metrics and uncertainty.
Do not infer human expert validation, deployment efficacy or language-model
roles from this synthetic observation. All such claims remain unestablished.
\end{document}
